\documentclass{article}
\usepackage{amsmath,amsthm,amssymb}
\usepackage{mathtext}
\usepackage[T2A]{fontenc}
\usepackage[utf8]{inputenc}
\usepackage[english]{babel}
\usepackage{graphicx}
\usepackage{hyperref}


\title{Cross-Era Generalization of Russian AI-Generated Text Detectors}
\author{Egor Sorvin}
\date{September 2026}


\begin{document}
\maketitle
\begin{abstract}
Detectors of machine-generated Russian text are trained on corpora built with 2020--2022 generators, while the
texts they meet today come from much newer models. We measure how such detectors transfer across corpora and
generator eras. Encoders fine-tuned on CoAT, zero-shot Binoculars with Russian-capable model pairs, and encoders
trained on corpus mixtures are evaluated on CoAT, LLMTrace, AINL-Eval 2025 and two new test sets generated with
T-pro~2.0 and Qwen3.8-27B. ruRoBERTa-large trained on CoAT drops from 0.95 AUROC on CoAT to 0.52--0.64 on
the newer corpora, and encoders trained on modern corpora score CoAT machine texts below human ones
(0.46--0.48 AUROC). With Russian-capable model pairs, Binoculars separates free continuations from human text
with up to 0.84 AUROC and paraphrases with at most 0.64; adding its features to a fine-tuned encoder changes
AUROC by less than 0.04. Mixtures that include LLMTrace reach 0.79--0.82 AUROC on the 2026 test sets, but the
training and test corpora share generator families, so the gain cannot be attributed to corpus diversity
alone. Code:
\url{https://github.com/egorsorvin/ru-ai-text-detection}.
\end{abstract}


\section{Introduction}
Detectors of AI-generated text are used in education, content moderation and fraud screening. For Russian, the
main public resources come from the RuATD 2022 shared task and its extension CoAT~\cite{shamardina2022ruatd,
shamardina2025coat}, built with generators such as ruGPT-3 and ruT5. Generators have changed since then: current
models write fluent text, and a detector that learned the artifacts of older models may miss it. In English,
detection is known to degrade on unseen generators and domains~\cite{dugan2024raid}, but this has not been
measured across Russian corpora of different eras.

We evaluate three kinds of detectors on three public corpora, CoAT, LLMTrace~\cite{tolstykh2025llmtrace} and
AINL-Eval~\cite{batura2025ainl}, under a leave-one-corpus-out protocol. We also build two test sets that pair
CoAT human texts with outputs of 2026 generators, so the human side stays fixed and only the generator changes.
Our contributions are:
\begin{itemize}
    \item a cross-corpus and cross-era evaluation of Russian detectors, including two new held-out test sets;
    \item evidence that the shift goes both ways: ruRoBERTa-large trained on CoAT flags 14--17\% of machine
    texts from 2026 generators, while encoders trained on modern corpora score ruGPT-3 outputs from CoAT as more
    human than human texts (0.32--0.38 AUROC);
    \item an analysis of Binoculars~\cite{hans2024binoculars} for Russian: with Russian-capable pairs it separates
    free continuation from human text with 0.69--0.84 AUROC and paraphrase or simplification with 0.51--0.66,
    and a threshold fitted on CoAT gives 0.50--0.63 balanced accuracy on newer corpora;
    \item a negative result: adding Binoculars features to a multi-corpus encoder changes AUROC by $-0.03$ to
    $+0.03$ in a single run, and the gains from corpus mixtures are confounded with generator families shared
    between training and test corpora.
\end{itemize}

\subsection{Team}
\textbf{Egor Sorvin} designed the study, built the test sets, ran the experiments and wrote the report.


\section{Related Work}
\label{sec:related}

\paragraph{Russian corpora.} RuATD 2022~\cite{shamardina2022ruatd} and its extension
CoAT~\cite{shamardina2025coat} contain texts from 13 generators of 2020--2021 fine-tuned for translation,
paraphrasing, summarization and simplification, plus free generation. LLMTrace~\cite{tolstykh2025llmtrace}
covers generators of 2023--2025, including the GPT-4 family, GigaChat, YandexGPT and Qwen.
AINL-Eval 2025~\cite{batura2025ainl} contains scientific abstracts by GPT-4-Turbo, Gemma-2-27B, Llama-3.3-70B and
withheld models. Each work evaluates detectors within its own corpus.

\paragraph{Generalization.} M4~\cite{wang2024m4}, MULTITuDE~\cite{macko2023multitude} and
RAID~\cite{dugan2024raid} show that trained detectors lose accuracy on unseen generators and domains.
Kirichenko~\cite{kirichenko2026sameorigin} shows that on Russian data detector scores also follow the text
source, so cross-corpus results mix a generator shift with a source shift.

\paragraph{Zero-shot detection.} DetectGPT~\cite{mitchell2023detectgpt}, Fast-DetectGPT~\cite{bao2024fastdetectgpt}
and Binoculars~\cite{hans2024binoculars} use the probabilities of a language model instead of training.
Binoculars detects over 90\% of ChatGPT texts at a 0.01\% false-positive rate in English, but LLMTrace reports an
AI-class F1 of 0.02 on Russian. Paraphrasing evades such detectors~\cite{krishna2023paraphrasing}.

Published results are in Table~\ref{tab:prior}. They are not directly comparable with ours: the official CoAT and
AINL-Eval test labels are closed.

\begin{table}[tbh!]
\begin{center}
\begin{tabular}{|l|l|l|c|}
\hline
Work & Detector & Metric & Value \\
\hline
RuATD 2022, binary & TF-IDF + logistic regression & Accuracy & 0.642 \\
 & BERT baseline & Accuracy & 0.797 \\
 & Best system, fine-tuned encoders & Accuracy & 0.830 \\
 & Human annotators & Accuracy & 0.667 \\
\hline
CoAT, binary & ruRoBERTa-large & Accuracy & 0.864 \\
 & Human annotators & Accuracy & 0.66 \\
\hline
LLMTrace, Russian & Fine-tuned Mistral-7B & AI-class F1 & 0.986 \\
 & Binoculars, zero-shot & AI-class F1 & 0.02 \\
\hline
\end{tabular}
\caption{Published results on the corpora used in this work, as reported by their authors.}
\label{tab:prior}
\end{center}
\end{table}


\section{Model Description}
% TODO


\section{Dataset}
% TODO


\section{Experiments}
\subsection{Metrics}
% TODO

\subsection{Experiment Setup}
% TODO

\subsection{Baselines}
% TODO


\section{Results}
% TODO


\section{Conclusion}
% TODO


\bibliographystyle{apalike}
\bibliography{lit}
\end{document}
